
Language models that perform well on capability benchmarks can simultaneously exhibit correlated failures in calibration, hallucination resistance, and adversarial robustness. This paper investigates whether these failure modes share a common latent dimension — termed \textit{epistemic reliability} — that is independent of raw MMLU performance. Using the YouRA multi-hypothesis evaluation framework, partial Spearman correlation, factor analysis, and leave-one-out prediction were applied to a synthetic score matrix conforming to a hypothesized latent structure, representing a population of N=30 open-weight instruction-tuned LLMs (7B–70B parameters, 8 model families). Under these synthetic conditions, a single epistemic reliability factor accounts for 72.1\% of shared variance across five trustworthiness metrics (KMO = 0.879). The calibration–hallucination partial correlation survives MMLU capability control with a survival fraction of 0.943 (partial ρ = −0.758, BCa 95\% CI [−0.894, −0.504]). A leave-one-out composite epistemic predictor achieves AUC = 0.739 for adversarial failure classification, but its incremental advantage over a capability-only baseline is statistically inconclusive (ΔAUC = 0.051, BCa 95\% CI [−0.194, 0.449]). All quantitative results are derived from a synthetic score matrix designed to conform to the hypothesized latent structure; they constitute pipeline validation under controlled conditions, not empirical measurements of real LLM behavior. Real-data replication via the pre-registered lm-evaluation-harness pipeline is the immediate scientific priority. The overall verdict across sub-hypotheses is PARTIALLY SUPPORTED.

